% Palestra de Biblioteconomia: letra grande, imagens e pouco texto
% Modelo do ÉguaTex (https://eguatex.github.io), em domínio público pela licença CC0 1.0:
% use, modifique e distribua sem pedir permissão e sem citar a origem.
%
% Compile com pdfLaTeX (o ÉguaTex faz isso sozinho).
%
% Numa palestra, o slide apoia a fala: uma ideia, uma imagem e no máximo uma frase.
% As imagens de exemplo (example-image-a, -b, -c) vêm do pacote mwe. Para usar as suas,
% coloque o arquivo na pasta do projeto e troque o nome em \includegraphics.

\documentclass[17pt,aspectratio=169]{beamer}

% ---------- idioma e fontes ----------
\usepackage[T1]{fontenc}
\usepackage[utf8]{inputenc}
\usepackage[brazilian]{babel}
\usepackage{lmodern}
\usepackage[default]{lato}

% ---------- visual limpo: sem barras, sem rodapé ----------
\definecolor{tinta}{RGB}{38,38,52}
\definecolor{destaque}{RGB}{214,82,54}
\setbeamercolor{normal text}{fg=tinta}
\setbeamercolor{frametitle}{fg=destaque}
\setbeamercolor{alerted text}{fg=destaque}
\setbeamerfont{frametitle}{size=\large, series=\bfseries}
\setbeamertemplate{navigation symbols}{}
\setbeamertemplate{footline}{%
  \hfill{\color{gray}\tiny\insertframenumber\hspace{6pt}}\vspace{4pt}}

% Slide de imagem cheia com uma frase por cima, numa faixa escura.
% Uso: \imagemcheia{arquivo-da-imagem}{Frase curta}
\newcommand{\imagemcheia}[2]{%
  \begingroup
  \setbeamertemplate{background}{\includegraphics[width=\paperwidth,height=\paperheight]{#1}}%
  \begin{frame}[plain]
    \vfill
    \begin{beamercolorbox}[wd=\paperwidth, sep=8pt]{faixa}
      \centering\color{white}\bfseries #2
    \end{beamercolorbox}
    \vspace{10pt}
  \end{frame}
  \endgroup}
\setbeamercolor{faixa}{bg=tinta}

% ---------- instruções de preenchimento ----------
% As notas em cinza explicam o que pôr em cada slide.
% Troque \instrucoestrue por \instrucoesfalse para esconder todas de uma vez.
\newif\ifinstrucoes
\instrucoestrue
\newcommand{\instrucao}[1]{\ifinstrucoes\par\vfill{\color{gray}\tiny\itshape #1\par}\fi}

% ---------- dados da capa: preencha aqui ----------
\title{A biblioteca como lugar de encontro}
\author{Nome do palestrante}
\institute{Universidade Federal Exemplo}
\date{Evento, Cidade, Ano}

\begin{document}

% ======================================================================
\begin{frame}[plain]
  \begin{columns}[c]
    \begin{column}{0.55\textwidth}
      {\usebeamercolor[fg]{frametitle}\Large\bfseries\inserttitle\par}
      \bigskip
      {\small\insertauthor\par}
      {\tiny\insertinstitute\par\insertdate\par}
    \end{column}
    \begin{column}{0.42\textwidth}
      \includegraphics[width=\linewidth]{example-image-1x1}
    \end{column}
  \end{columns}
\end{frame}

% ======================================================================
\begin{frame}
  \centering\vfill
  Quem entrou numa biblioteca\\na última semana?
  \vfill
  \instrucao{Abra com uma pergunta para a plateia. Espere as mãos levantadas antes de passar o slide.}
\end{frame}

% ======================================================================
\imagemcheia{example-image-a}{Leitura também é convivência}

% ======================================================================
\begin{frame}{Três ideias}
  \begin{columns}[T]
    \begin{column}{0.31\textwidth}
      \centering
      \includegraphics[width=\linewidth]{example-image-a}\\[4pt]
      \small Acolher
    \end{column}
    \begin{column}{0.31\textwidth}
      \centering
      \includegraphics[width=\linewidth]{example-image-b}\\[4pt]
      \small Mediar
    \end{column}
    \begin{column}{0.31\textwidth}
      \centering
      \includegraphics[width=\linewidth]{example-image-c}\\[4pt]
      \small Conectar
    \end{column}
  \end{columns}
  \instrucao{Resuma a palestra em três palavras com uma imagem cada. Fale um minuto sobre cada uma.}
\end{frame}

% ======================================================================
\begin{frame}
  \centering\vfill
  {\Huge\bfseries\alert{1 em 3}}\par\medskip
  \small visitantes vêm só para estudar\par
  {\tiny\color{gray} (número de exemplo: troque pelo seu dado e pela fonte)}
  \vfill
  \instrucao{Um número grande por slide fica na memória. Diga de onde ele veio.}
\end{frame}

% ======================================================================
\begin{frame}
  \vfill
  \begin{quote}
    \itshape ``A biblioteca é um organismo em crescimento.''
  \end{quote}
  \hfill{\small Ranganathan, 1931}
  \vfill
  \instrucao{Uma citação curta, com autor e ano, ajuda a fechar um bloco da fala.}
\end{frame}

% ======================================================================
\imagemcheia{example-image-16x9}{O que você leva daqui?}

% ======================================================================
\begin{frame}[plain]
  \centering\vfill
  {\Large\bfseries Obrigado!}\par\bigskip
  \small Nome do palestrante\par
  \tiny email@exemplo.com
  \vfill
\end{frame}

\end{document}
